\chapter{Parametri e configurabilità}
\label{ch:param}

\section{Parametri di build (synthesis-time)}
L'architettura hardware è fissata alla sintesi tramite i parametri Verilog seguenti.
Determinano il datapath contenuto nel bitstream e il suo \emph{soffitto} di capacità.

\begin{tabularx}{\textwidth}{L{3.0cm} C{1.8cm} Y}
\toprule
\rowh \thd{Parametro} & \thd{Default} & \thd{Significato} \\
\midrule
\code{DATA\_WIDTH}   & 8   & Larghezza dei dati (INT8). \\
\rowa \code{ACC\_WIDTH} & 32 & Larghezza dell'accumulatore (INT32). \\
\code{N\_INPUTS}     & 32 / 256 & Numero massimo di ingressi per neurone (baseline benchmark: 256). \\
\rowa \code{N\_NEURONS} & 1 / 4 & Numero massimo di neuroni per layer. \\
\code{PARALLEL}      & 8   & MAC hardware simultanei per neurone; deve dividere \code{N\_INPUTS} e conviene sia potenza di due. \\
\rowa \code{N\_LAYERS} & 4 & Numero massimo di layer concatenabili da \code{layer\_sequencer}. \\
\code{ADDR\_WIDTH}   & 23  & Larghezza dell'indirizzo di byte (8~MB). \\
\rowa \code{MEM\_DATA\_WIDTH} & 16 & Larghezza del bus dati fisico PSRAM. \\
\code{CLK\_FREQ\_MHZ} & 80 & Frequenza usata per le formule di temporizzazione PSRAM (va allineata all'oscillatore reale). \\
\bottomrule
\end{tabularx}

\begin{fnwarn}[Vincolo \texttt{N\_INPUTS} \% \texttt{PARALLEL}]
\code{PARALLEL} deve dividere esattamente \code{N\_INPUTS}, altrimenti scatta il guard
di elaborazione (§\ref{ch:datapath}). Lo stesso vincolo vale a runtime su
\code{n\_inputs\_real}.
\end{fnwarn}

\section{Larghezza di rete a runtime}
Un singolo bitstream serve qualunque topologia \emph{fino} al massimo di build. La
larghezza reale di ciascuna esecuzione è un valore separato, impostato dall'host:
\begin{itemize}
\item \code{n\_inputs\_real} --- ingressi realmente usati in questa esecuzione (deve
essere multiplo di \code{PARALLEL});
\item \code{n\_neurons\_real} --- neuroni realmente calcolati in questa esecuzione.
\end{itemize}
Entrambi hanno default pari al massimo di build, così ogni chiamante che li lascia
scollegati elabora l'intera larghezza come prima dell'introduzione delle porte.

\begin{fnnote}[Terminazione anticipata reale]
Non si tratta di semplice contabilità di indirizzi: i due valori limitano
direttamente i loop hardware (letture X/W di \code{neuron\_memory}, conteggio gruppi
MAC di \code{neuron\_parallel} e lunghezza della copia ping-pong per \code{RUN\_NETWORK}).
Un layer più stretto \emph{calcola} e \emph{copia} davvero più in fretta e non richiede
zero-padding della RAM per la coda non usata: i dati oltre
\code{n\_inputs\_real}/\code{n\_neurons\_real} non vengono mai letti.
\end{fnnote}

Questo permette a una rete di rastremarsi dentro una sola esecuzione concatenata, ad
esempio $256\to64\to16\to4$, con ogni layer che dichiara la propria larghezza reale
nella tabella descrittori (cap.~\ref{ch:seq}).

\subsection{Risparmio misurato}
La terminazione anticipata è stata misurata end-to-end:
\begin{tabularx}{\textwidth}{L{5.5cm} C{3.0cm} Y}
\toprule
\rowh \thd{Test} & \thd{Cicli} & \thd{Confronto} \\
\midrule
\code{neuron\_parallel\_tb.v} (T7) & 3 vs 6 & ridotto vs pieno, con dati ``spazzatura'' nelle corsie saltate (prova che non vengono lette). \\
\rowa \code{neuron\_memory\_tb.v} (T5) & 209 vs 788 & 8-di-32 vs 32 pieni, attraverso lo stack PSRAM reale. \\
\bottomrule
\end{tabularx}

\section{Configurazioni caratterizzate}
Alcune combinazioni convalidate in simulazione e/o sintesi:
\begin{tabularx}{\textwidth}{C{2.0cm} C{2.0cm} C{2.0cm} Y}
\toprule
\rowh \thd{N\_INPUTS} & \thd{N\_NEURONS} & \thd{PARALLEL} & \thd{Note} \\
\midrule
32 & 4 & 8 & Primo test parametrico funzionale (Fase~1). \\
\rowa 256 & 4 & 2/4/8/16 & Sweep di benchmark del datapath (Fase~7). \\
32 & 1..3 & 8 & Integrazione memoria mono/multi-neurone (Fase~3). \\
\rowa 4 & 4 & 2 & Test end-to-end \code{RUN\_NETWORK} a 2 layer su SPI reale. \\
\bottomrule
\end{tabularx}

\section{Riepilogo build contro runtime}
\begin{center}
\begin{tikzpicture}[font=\footnotesize,node distance=6mm]
 \node[fnblockD,minimum width=54mm,minimum height=15mm](b){\textbf{BUILD (sintesi)}\\[2pt]
   {\scriptsize N\_INPUTS, N\_NEURONS, N\_LAYERS,}\\{\scriptsize PARALLEL, DATA\_WIDTH, ACC\_WIDTH}\\{\scriptsize $\Rightarrow$ soffitto della macchina}};
 \node[fnblockT,right=16mm of b,minimum width=54mm,minimum height=15mm](r){\textbf{RUNTIME (host, SPI)}\\[2pt]
   {\scriptsize n\_inputs\_real, n\_neurons\_real,}\\{\scriptsize attivazione, num\_layers, pesi/bias}\\{\scriptsize $\Rightarrow$ rete effettiva}};
 \draw[fnbus] (b) -- node[fnlbl,above]{$\le$} (r);
\end{tikzpicture}
\end{center}
